\chapter{APPLICATIONS}


% \section{Markov Chains}
\textit{October 14th.}

\begin{example}
    A monkey randomly chooses an alphabet key on a typewriter uniformly; let $u_{i}$ be the key chosen at $i$-th time. What is the chance the monkey types out ``ISP IS BORING''?
\end{example}
In particular, let
\begin{align}
    \tau = \min\{n \geq 11 \mid (u_{n-10},\ldots,u_{n}) = (I,S,P,I,S,B,O,R,I,N,G)\}.
\end{align}
We ask if $\E[\tau]$ is finite. The Markov chain method is as follows; for all $k \geq 1$, define $Y_{k} = (u_{k},\ldots,u_{k+10})$. As $u_{i} \sim \Unif\{A,B,\ldots,Z\}$, $\{Y_{k}\}_{k \geq 1}$ turns into a Markov chain. Using $h(x) = \E[\tau \mid Y_{1} = x]$, we may set up a recurrence relation and solve for $h$ by hand to obtain an answer. This is computationally expensive by hand.

A better approach is to use martingales. Let $(Y_{1},\ldots,Y_{11}) = (I,S,P,\ldots,I,N,G)$, and set $Z_{0}^{1} = 1$ to be the capital of gambler $1$ at time $t = 0$. At time $t = 1$, the gambler places a bet that $u_{1} = I$. If $u_{1} = I$, then the gambler gets $26$ capital, and if $u_{1} \neq I$, the gambler gets $0$ capital and exits the casino. The expected gain here is then
\begin{align}
    \frac{1}{26}(26-1) + \frac{25}{26}(-1) = 0.
\end{align}
At time $t = 2$, the gambler bets all the gain, $26$ capital, that $u_{2} = S$. If $u_{2} = S$, the gambler gets $(26)^{2}$ capital, and if $u_{2} \neq S$, the gambler loses and exits the casino. The expected gain here is, again, $0$. Repeat this until $t = 11$, and $Z_{t}^{1} = Z_{11}^{1}$ for all $t \geq 10$, where $Z_{n}^{1}$ is the current capital; that is,
\begin{align}
    Z_{n}^{1} = \begin{cases}
        26 \cdot Z_{n-1}^{1} \1_{\{u_{n} = \psi_{n}\}} &\text{ if } 1 \leq n \leq 11,\\
        Z_{11}^{1} &\text{ if } n \geq 12.
    \end{cases}
\end{align}
Now, a second gambler enters the casino at time $t = 2$ and places a bet of $1$ that $u_{2} = I$. If $u_{2} = I$, the second gambler gets $26$ capital, and $0$ otherwise. Continue similarly till $t = 12$. So, for any $k \geq 1$, gambler $k$ enters at time $t = k$ and starts with capital $1$ for bet $u_{k} = I$, with capital at each step as
\begin{align}
    Z_{n}^{k} = \begin{cases}
        1 &\text{ if } n < k,\\
        26 \cdot Z_{n-1}^{k} \1_{\{u_{n}=\psi_{n-k+1}\}} &\text{ if } k \leq n \leq k+10,\\
        Z_{k+10}^{k} &\text{ if } n \geq k + 11.
    \end{cases}
\end{align}
Now, finally, define
\begin{align}
    M_{n} = \sum_{k=1}^{n}(1-Z_{n}^{k}) = n - \sum_{k=1}^{n}Z_{n}^{k}.
\end{align}
the gain of the casino. We make three claims:
\begin{enumerate}
    \item $(M_{n})_{n \geq 1}$ is a martingale.
    \item $\tau$ is a stopping time for this martingale.
    \item The OST applies.
\end{enumerate}
To show $(M_{n})_{n \geq 1}$ is a martingale, we first have
\begin{align}
    \abs{M_{n}} \leq n + n(26)^{11} < \infty,
\end{align}
so they are integrable. Moreover,
\begin{align}
    \E[M_{n} \mid M_{n-1},\ldots,M_{1}] &= n - \E\left[ \sum_{k=1}^{n-11} Z_{n-1}^{k} + \sum_{k=n-10}^{n}26 \cdot Z_{n-1}^{k} \1_{\{u_{n}=\psi_{n-k+1}\}} \mid u_{n-1},\ldots,u_{1} \right] \notag \\
    &= n - \sum_{k=1}^{n-11} Z_{n-1}^{k} + 26\sum_{k=n-10}^{n} Z_{n-1}^{k} \Prob(u_{n} = \psi_{n-k+1}) \notag \\
    &= n - \sum_{k=1}^{n-11} Z_{n-1}^{k} - \sum_{k=n-10}^{n} Z_{n-1}^{k} = n-1-\sum_{k=1}^{n-1}Z_{n-1}^{k} = M_{n-1},
\end{align}
showing that $(M_{n})_{n \geq 1}$ is a martingale. For the second claim, note that $\{\tau \leq n\} = \bigcup_{k=1}^{n}\{(u_{k},\ldots,u_{k+10}) = (I,S,P,\ldots,I,N,G)\} \in \sigma(u_{1},\ldots,u_{n})$, so $\tau$ is a stopping time. Now, $M_{n} = n-\sum_{k=1}^{n} Z_{n}^{k}$, and $\E[M_{1}] = 1 - \E[Z_{1}^{1}] = 0$. Note that $Z_{\tau}^{k} = 0$ for all $k \leq \tau-10$, and $Z_{\tau}^{k} = 0$ for all $\tau-10 < k \leq \tau$ except for $k = \tau-10$, where $Z_{\tau}^{\tau-10} = 26^{11}$. If the OST applies, then $\E[M_{\tau}] = \E[M_{1}] = 0$, so
\begin{align}
    \E[\tau] = \E[M_{\tau}] + (26)^{11} = (26)^{11} < \infty.
\end{align}
Thus, we are only left to show that the OST applies. If we let $W_{k}$ denote the $k$-th 11-block of $u$'s, then the $W_{k}$'s are independent and identically distributed. If $S$ is the first $k$ such that $W_{k} = (I,S,P,\ldots,I,N,G)$, then $\tau \leq 11S$. But this implies
\begin{align}
    \E[\tau] \leq 11\E[S] = 11 \cdot 26^{11} < \infty,
\end{align}
so $\tau$ is integrable. Now, the increments in $M_{n}$ are handled as
\begin{align}
    \abs{M_{n}-M_{n-1}} &= \abs{\sum_{k=1}^{n-1}(Z_{n-1}^{k}-Z_{n}^{k})-Z_{n}^{n}+1} \leq \abs{\sum_{k=1}^{n-1}(Z_{n-1}^{k}-Z_{n}^{k})} + 25 \notag \\ &\leq 25\sum_{k=n-10}^{n-1}\abs{Z_{n-1}^{k}} + 25 \leq 25(26)^{11}10+25.
\end{align}
The OST applies, and we are done.

\begin{example}
    Let $S_{n}$ be the permutation group on $n$ elements, $S_{n} = \{\sigma: \{1,\ldots,n\} \to \{1,\ldots,n\} \mid \sigma \text{ is bijective}\}$. Let $\pi \in S_{n}$. An increasing subsequence of $\pi$ is a sequence $1 \leq i_{1} < i_{2} < \ldots < i_{k} \leq n$ such that $\pi(i_{1}) < \pi(i_{2}) < \ldots < \pi(i_{k})$. Let $L_{n}(\pi)$ denote the length of the longest increasing subsequence of $\pi$, where $\pi$ is chosen uniformly at random from $S_{n}$. What is the asymptotic behavior of $L_{n} = L_{n}(\pi)$ as $n \to \infty$?
\end{example}

We show the following theorem.
\begin{theorem}
    The following hold for $L_{n}$.
    \begin{enumerate}
        \item $\E[L_{n}] \geq \sqrt{n}$.
        \item $\limsup_{n \to \infty} \frac{\E[L_{n}]}{\sqrt{n}} \leq e$.
        \item $\Prob(\abs{L_{n}-\E[L_{n}]} \geq t\sqrt{n}) \leq 2e^{-t^{2}/2}$.
    \end{enumerate}
\end{theorem}
\noindent From these statements, one can conclude that
\begin{align}
    1 \leq \liminf_{n \to \infty} \frac{\E[L_{n}]}{\sqrt{n}} \leq \limsup_{n \to \infty} \frac{\E[L_{n}]}{\sqrt{n}} \leq e,
\end{align}
and
\begin{align}
    \frac{L_{n}}{\sqrt{n}} - \frac{\E[L_{n}]}{\sqrt{n}} \toP 0.
\end{align}

\noindent We make use of the following fact, that if $x_{1},\ldots,x_{n}$ are real numbers $L(x_{1},\ldots,x_{n})$ and $D(x_{1},\ldots,x_{n})$ denoting the lenghts of the longest increasing and decreasing subsequences of $x_{1},\ldots,x_{n}$, then
\begin{align}
    L(x_{1},\ldots,x_{n})D(x_{1},\ldots,x_{n}) \geq n.
\end{align}
Note that $D_{n}(\pi)$, the length of the longest decreasing subsequence of $\pi$, is equal in distribution to $L_{n}(\pi)$, so $\E[L_{n}]^{2} \geq n$, giving the first statement. For the second statement, let $N_{k}$ denote the number of increasing subsequences of length $k$. Then
\begin{align}
    N_{k} = \sum_{1 \leq i_{1} < \ldots < i_{k} \leq n} \1_{\{\pi(i_{1}) < \ldots < \pi(i_{k})\}}.
\end{align}
Then $\E[N_{k}] = \binom{n}{k}\frac{1}{k!}$. Then using Markov's inequality and Stirling's formula gives
\begin{align}
    \Prob(L_{n} \geq k) = \Prob(N_{k} \geq 1) \leq \E[N_{k}] = \binom{n}{k} \frac{1}{k!} \implies \E[L_{n}] = \sum_{k=1}^{n} \Prob(L_{n} \geq k) \leq \sum_{k=1}^{\infty} \binom{n}{k} \frac{1}{k!} \leq \sqrt{n}e.
\end{align}
For the last statement, let $\cA_{k} = \sigma(\pi(1),\ldots,\pi(k))$, and set $M_{k} = \E[L_{n} \mid \cA_{k}]$ for $n \geq 1$, and $M_{0} = \E[L_{n}]$. Then $(M_{k})_{0 \leq k \leq n}$ is a martingale (show this). Let $D_{k} = M_{k}-M_{k-1}$. Then
\begin{align}
    \E[D_{k} \mid \cA_{k-1}] = \E[M_{k} \mid \cA_{k-1}] - M_{k-1} = 0.
\end{align}
Also, there exists $c_{k}$ such that $\abs{D_{k}} \leq c_{k}$ (show this). This implies that
\begin{align}
    \E[e^{\lambda D_{k}} \mid \cA_{k-1}] \leq e^{\lambda^{2}c_{k}^{2}/2}
\end{align}
or
\begin{align}
    \E[e^{\lambda (M_{n}-M_{0})}] = \E[e^{\lambda \sum_{k=1}^{n} D_{k}}] = \E[e^{\sum_{k=1}^{n-1} D_{k}} \E[e^{\lambda D_{n}} \mid \cA_{n-1}]] \leq e^{\lambda^{2}\sum_{k=1}^{n} c_{k}^{2}/2}.
\end{align}
Using this, we have
\begin{align}
    \Prob(M_{n}-M_{0} \geq t) \leq \Prob(e^{\lambda(M_{n}-M_{0})} \geq e^{\lambda t}) \leq e^{-\lambda t} \E[e^{\lambda(M_{n}-M_{0})}] \leq e^{-\lambda t + \lambda^{2}\sum_{k=1}^{n} c_{k}^{2}/2}.
\end{align}
One can show that $c_{k} = 1$ works, so that
\begin{align}
    \Prob(M_{n}-M_{0} \geq t) \leq e^{-\lambda t + n\lambda^{2}/2}.
\end{align}
Minimizing the right-hand side over $\lambda$ gives $\lambda = t/n$, so
\begin{align}
    \Prob(M_{n}-M_{0} \geq t) \leq e^{-t^{2}/2n}.
\end{align}
This gives the desired result, since $\sqrt{n}(M_{n}-M_{0}) = L_{n}-\E[L_{n}]$.
\\

\textit{October 21st.}

\begin{example}
    The Erdos-Renyi random graph $G(n,p)$ is a random graph on $n$ vertices where each edge is included independently with probability $p$. We show results on this $G(n,p)$.
\end{example}
Let $p = \lambda/n$ for some $\lambda > 0$. Then note that
\begin{align}
    \Prob(\deg(i) = k) = \binom{n-1}{k}p^{k}(1-p)^{n-1-k} \to \frac{\lambda^{k}}{k!}e^{-\lambda},
\end{align}
as $n \to \infty$. We ask if there is a threshold of $\lambda$ that corresponds to connectivity of $G(n,\lambda/n)$. The following theorem applies.

\begin{theorem}
    Let $G(n,p)$ be the Erdos-Renyi random graph with $p = \lambda/n$. Then
    \begin{enumerate}
        \item If $\lambda - \log n \to -\infty$ as $n \to \infty$, then $\Prob(G(n,p) \text{ is connected}) \to 0$.
        \item If $\lambda - \log n \to +\infty$ as $n \to \infty$, then $\Prob(G(n,p) \text{ is connected}) \to 1$.
        \item If $\lambda - \log n \to t$ as $n \to \infty$, then $\Prob(G(n,p) \text{ is connected}) \to e^{-e^{-t}}$.
    \end{enumerate}
\end{theorem}

\begin{proof}
    Let $Y_n$ be the number of isolated vertices. Then
    \begin{align}
        \E[Y_n]=n(1-p)^{n-1},
        \qquad
        \E[Y_n(Y_n-1)]=n(n-1)(1-p)^{2n-3}.
    \end{align}
    Consequently, Chebyshev's inequality gives
    \begin{align}
        \Prob(Y_n=0)\leq \frac{1}{\E[Y_n]}+\frac{\lambda}{n-\lambda}.
    \end{align}
    (The second term may be replaced by $(\lambda-1)/(n-\lambda)$.)

    If $\lambda-\log n\to-\infty$, then $\E[Y_n]\to\infty$, and
    the above implies $\Prob(Y_n=0)\to0$. Since a connected
    graph has no isolated vertex, the first assertion follows.

    We next show that, whenever $\lambda\geq \frac12\log n$,
    \begin{align}
        \Prob(G(n,p)\text{ is disconnected and }Y_n=0)\longrightarrow0.
        \label{eq:no-small-components}
    \end{align}
    For $2\leq k\leq n/2$, let $X_k$ be the number of components of size $k$.
    A component counted by $X_k$ contains a spanning tree, and hence
    \begin{align}
        \E[X_k]
        \leq \binom{n}{k}k^{k-2}p^{k-1}(1-p)^{k(n-k)}.
        \label{eq:component-bound}
    \end{align}
    Summing this bound and using $p=\lambda/n$ shows
    \begin{align}
        \sum_{k=2}^{\lfloor n/2\rfloor}\E[X_k]=o(1).
        \label{eq:component-sum}
    \end{align}
    Indeed, use $(1-p)^{k(n-k)}\leq e^{-\lambda k(1-k/n)}$ in
    \eqref{eq:component-bound}, and split the sum at $k=\sqrt n$; the first
    part is bounded by a geometric series with ratio $O(\lambda/n)$, while
    the second part is $O(n^2e^{-c\lambda\sqrt n})$ for an absolute $c>0$.
    Thus Markov's inequality proves \eqref{eq:no-small-components}, since a
    disconnected graph with no isolated vertices has a component of size
    between $2$ and $n/2$.

    If $\lambda-\log n\to+\infty$, then
    \begin{align}
        \Prob(Y_n\geq1)\leq \E[Y_n]\leq ne^{-\lambda+\lambda/n}\to0.
    \end{align}
    Combining this with \eqref{eq:no-small-components} proves the second
    assertion.

    Finally, suppose $\lambda-\log n\to t$. For every fixed $r$,
    \begin{align}
        \E[(Y_n)_r]
        =(n)_r(1-p)^{r(n-r)+\binom r2}\longrightarrow e^{-rt}.
    \end{align}
    Hence $Y_n$ converges in distribution to a Poisson random variable with
    mean $e^{-t}$, and therefore
    \begin{align}
        \Prob(Y_n=0)\longrightarrow e^{-e^{-t}}.
    \end{align}
    Since $\lambda\geq\frac12\log n$ eventually, \eqref{eq:no-small-components}
    gives
    \begin{align}
        \Prob(G(n,p)\text{ is connected})
        =\Prob(Y_n=0)+o(1)\longrightarrow e^{-e^{-t}}.
    \end{align}
    This completes the proof.
\end{proof}

\textit{October 28th.}

\begin{example}
    Consider the random walk where $\Prob(X=0) = \frac{1}{2} = \Prob(X=1)$, and $S_{n} = \sum_{i=1}^{n}X_{i}$. Strong law of large numbers gives $S_{n}/n \to 1/2$ almost surely, and the central limit theorem gives $(S_{n}-n/2)/\sqrt{n} \xrightarrow{d} N(0,1/4)$.
\end{example}

Given $a > \frac{1}{2}$, we ask for the asymptotic behavior of $\Prob(S_{n} \geq an)$.

\begin{theorem}
    As $n \to \infty$,
    \begin{align}
        \frac{1}{n} \log \Prob(S_{n} \geq an) \to -I(a)
    \end{align}
    where
    \begin{align}
        I(z) = \begin{cases}
            \log 2 + z \log z + (1-z)\log(1-z) &\text{ if } z \in [0,1],\\
            \infty &\text{ otherwise.}
        \end{cases}
    \end{align}
\end{theorem}

The function $I$ above is called the \eax{rate function} for the random walk.

\begin{proof}
    Note that for $a \in (1,\infty)$, the proof is trivial. For $a \in (1/2,1]$, we have $\Prob(S_{n} \geq an) = \frac{1}{2^{n}} \sum_{k \geq an} \binom{n}{k}$. Letting $Q_{n}(a) = \max_{k \geq an} \binom{n}{k}$, we have
    \begin{align}
        \frac{1}{2^{n}} Q_{n}(a) \leq \Prob(S_{n} \geq an) \leq \frac{n+1}{2} Q_{n}(a).
    \end{align}
    Using Stirling's formula, we can approximate $Q_{n}(a)$ as
    \begin{align}
        Q_{n}(a) \sim \frac{1}{\sqrt{2\pi an(1-a)n}} \frac{n^{n}}{(an)^{an}((1-a)n)^{(1-a)n}} = \frac{1}{\sqrt{2\pi an(1-a)n}} \left(\frac{1}{a^{a}(1-a)^{(1-a)}}\right)^{n}.
    \end{align}
    From here, we have
    \begin{align}
        \frac{1}{n} \log Q_{n}(a) \to -a\log a - (1-a)\log(1-a),
    \end{align}
    and therefore
    \begin{align}
        \frac{1}{n} \log \Prob(S_{n} \geq an) \to -a\log a - (1-a)\log(1-a)
    \end{align}
    via the squeeze theorem.
\end{proof}

\begin{remark}
    Via switching of a few terms, one can show that if $a < 1/2$, then
    \begin{align}
        \frac{1}{n} \log \Prob(S_{n} \leq an) \to -I(1-a).
    \end{align}
\end{remark}

\begin{theorem}[\eax{Cramer's theorem}]
    Suppose $X_{1},X_{2},\ldots$ are independent and identically distributed with $X$ having moment generating function $\varphi(t) = \E[e^{tX}] < \infty$ for all $t \in \R$. Then for any $a > \E[X]$,
    \begin{align}
        \lim_{n \to \infty} \frac{1}{n} \log \Prob(S_{n} \geq an) = -I(a)
    \end{align}
    where the rate function $I$ is given by $I(z) = \sup_{t \in \R}\{tz - \log \varphi(t)\}$.
\end{theorem}

\begin{proof}
    Without the loss of generality, we may assume $\E[X] < 0$ and $a = 0$. Thus, we are equivalently show that $\frac{1}{n} \log \Prob(S_{n} \geq 0) \to \log \rho$ where $\rho = \inf_{t \in \R} \varphi(t)$. Note that $\varphi$ is convex and smooth with $\varphi(0) = 1$ and $\varphi'(0) = \E[X] < 0$. We do casework. If $\Prob(X < 0) = 1$, then $\Prob(S_{n} \geq 0) = 0$ for all $n$, and the result is trivial. If $\Prob(X \leq 0) = 1$ and $\Prob(X = 0) > 0$, then $\varphi$ is strictly decreasing, and $\rho = \lim_{t \to \infty} \varphi(t) = \Prob(X = 0)$. Then $\Prob(S_{n} \geq 0) = \Prob(X = 0)^{n}$, and the result follows. The final case is when $\Prob(X > 0) > 0$ and $\Prob(X < 0) > 0$. Here, with $\alpha > 0$,
    \begin{align}
        \Prob(S_{n} \geq 0) = \Prob(e^{\alpha S_{n}} \geq 1) \leq \E[e^{\alpha S_{n}}] = \varphi(\alpha)^{n}.
    \end{align}
    Since $\lim_{t \to \pm \infty} \varphi(t) = \infty$, there exists $\alpha > 0$ such that $\varphi(\alpha) = \rho$. With this $\alpha$,
    \begin{align}
        \limsup_{n \to \infty} \frac{1}{n} \log \Prob(S_{n} \geq 0) \leq \log \varphi(\alpha) = \log \rho.
    \end{align}
    We now utilize tilting. With $F$ the distribution function of $X$, define $\{\tilde{X}_{i}\}_{i \geq 1}$, a sequence of independent and identically distributed random variables with distribution function $\tilde{F}$ given by
    \begin{align}
        \tilde{F}(x) = \frac{1}{\rho} \int_{-\infty}^{x} e^{\alpha y} dF(y).
    \end{align}
    Then, $\E[\tilde{X}] = \frac{1}{\rho} \int_{-\infty}^{\infty} y e^{\alpha y} dF(y) = \frac{1}{\rho} \frac{d}{d\alpha} \int_{-\infty}^{\infty} e^{\alpha y} dF(y) = \frac{1}{\rho} \frac{d}{d\alpha} \varphi(\alpha) = 0$, and $\Var(\tilde{X}) = \tilde{\sigma}^{2} < \infty$. Let $\tilde{S}_{n} = \sum_{i=1}^{n} \tilde{X}_{i}$. Then, we have
    \begin{align}
        \Prob(\tilde{S}_{n} \geq 0) = \frac{1}{\rho^{n}} \E[e^{-\alpha \tilde{S}_{n}} \1_{\{\tilde{S}_{n} \geq 0\}}] \geq \frac{1}{\rho^{n}} e^{-\alpha C \tilde{\sigma} \sqrt{n}} \E[1_{\{\frac{\tilde{S}_{n}}{\sqrt{n} \tilde{\sigma}} \in [0,C]\}}].
    \end{align}
    Via the central limit theorem, the rightmost expectation converges to $\Prob(Z \in [0,C])$ where $Z \sim N(0,1)$. Choose $C$ such that $\Prob(Z \in [0,C]) > \frac{1}{4}$. This gives
    \begin{align}
        \liminf_{n \to \infty} \frac{1}{n} \log \Prob(\tilde{S}_{n} \geq 0) \geq \log \rho.
    \end{align}
    Combining the upper and lower bounds gives the desired result.
\end{proof}


\begin{example}
    \textit{Random walker in mobile Poisson traps}; consider a random walker starting at the origin. At every $0 \neq z \in \Z$, a $\Poi(\gamma)$ number of traps are placed. $\{X_{k}\}_{k \geq 1}$ denotes the simple symmetric random walk executed by the walker, and, for $y \in \Z$, $\{Y^{j,y}\}_{1 \leq j \leq N_{y}}$ denotes the simple symmetric random walk executed by the $j$-th trap at $y$, where $N_{y} \sim \Poi(\gamma)$ is the number of traps at $y$. When a trap meets the walker, the walker is killed with probability $q$.
\end{example}
We ask for the survival probability of the walker up to time $n$. We use $\xi(k,x)$ to denote the number of traps at site $k$ at time $x$. Thus, at time $k$, the walker is at $X_{k}$, and the number of traps at that site is $\xi(k,X_{k})$; the probability that the walker survives at time $k$ is $(1-q)^{\xi(k,X_{k})}$. Thus, the survival probability of the walker up to time $n$, for a fixed trajectory of the walker, is
\begin{align}
    \prod_{i=1}^{n} (1-q)^{\xi(i,X_{i})} = (1-q)^{\sum_{i=1}^{n} \xi(i,X_{i})}.
\end{align}
For a \textit{fixed} (random) trajectory $X$ of the walker, we denote
\begin{align}
    \sigma^{X}(n) = \E^{\xi}\left[ \prod_{i=1}^{n} (1-q)^{\xi(i,X_{i})} \right]
\end{align}
where the expectation is over the randomness of the traps. \eax{Pascal's principle} says that the survival probability is maximized when the walker does not move; that is,
\begin{align}
    \sigma^{X}(n) \leq \sigma^{\bv{0}}(n) = \E^{\xi}\left[ \prod_{i=1}^{n} (1-q)^{\xi(i,0)} \right].
\end{align}
To this end, we have
\begin{align}
    \sigma^{X}(n) &= \E^{\xi}\left[ (1-q)^{\sum_{i=0}^{n}\sum_{y \in Z} \sum_{j=1}^{N_{y}} \1_{Y_{i}^{j,y} = X_{i}}} \right] = \prod_{y \in \Z} \E^{Y}_{y} \E^{N_{y}} \left[ \prod_{1 \leq j \leq N_{y}} (1-q)^{\sum_{i=0}^{n} \1_{Y_{i}^{j,y} = X_{i}}} \right] \notag \\
    &= \prod_{y \in \Z} \sum_{k=0}^{\infty} e^{-\gamma} \frac{\gamma^{k}}{k!} \E^{Y}_{y} \left[ \prod_{1 \leq j \leq k} (1-q)^{\sum_{i=0}^{n} \1_{Y_{i}^{j,y} = X_{i}}} \right] = \prod_{y \in \Z} \sum_{k=0}^{\infty} e^{-\gamma} \frac{\gamma^{k}}{k!} \left( \E^{Y}_{y} \left[ (1-q)^{\sum_{i=0}^{n} \1_{Y_{i} = X_{i}}} \right] \right)^{k} \notag \\
    &= \prod_{y \in \Z} e^{-\gamma (1 - \E^{Y}_{y} \left[ (1-q)^{\sum_{i=0}^{n} \1_{Y_{i} = X_{i}}} \right])} = e^{-\gamma \sum_{y \in \Z} (1 - \E^{Y}_{y} \left[ (1-q)^{\sum_{i=0}^{n} \1_{Y_{i} = X_{i}}} \right])}.
\end{align}
Denote $W^{X}(n,y) = 1 - \E^{Y}_{y} \left[ (1-q)^{\sum_{i=0}^{n} \1_{Y_{i} = X_{i}}} \right]$. Thus, showing Pascal's principle amounts to showing that
\begin{align}
    \sum_{y \in \Z} W^{X}(n,y) \geq \sum_{y \in \Z} W^{\bv{0}}(n,y).
\end{align}
The proof is probabilistic in nature. 